\section{How to Choose the Source of the Small $\mathbf{Z}$ Matrix}
\label{sec:choosing-z}

The matrix $\mathbf{Z} \in \mathbb{R}^{n \times k}$ is the only free input of
DPCR (Section~\ref{sec:dpcr}): the projector $\tilde{\mathbf{P}}$, the cost
of each iteration and the iteration count all follow from it. Any full-rank
$\mathbf{Z}$ is admissible, because \eqref{eq:dpcr-interlacing} keeps the
deflated spectrum inside that of $\mathbf{P}^{-1}\mathbf{A}$, so we are free
to choose $\mathbf{Z}$ by how cheap it is to build and apply, not only by
how well it deflates. The ideal choice is the span of the eigenvectors of
$\mathbf{P}^{-1}\mathbf{A}$ for its $k$ smallest eigenvalues, which reduces
the effective condition number to
$\lambda_n(\mathbf{P}^{-1}\mathbf{A})/\lambda_{k+1}(\mathbf{P}^{-1}\mathbf{A})$,
but computing them costs more than the solve itself. Because the iteration
count also depends on the right-hand side, only the soft modes that
$\tilde{\mathbf{P}}\mathbf{b}$ actually excites need to be captured.
Practical sources therefore approximate this ideal from information that is
available before the solve: earlier solves in the sequence, the Krylov
vectors they produced, or the structure of $\mathbf{A}$.
Table~\ref{tab:z-sources-tried} lists the sources we have tested.

\paragraph{Test problems.}
All experiments use linear systems recorded from the ostrich engine while it
simulates Helhest, a three-wheeled robot. Its chassis and three wheels are
four rigid bodies, connected by revolute joints, and every wheel is driven by
a velocity controller. The robot drives at constant wheel speed in two
scenes:
\begin{itemize}
  \item \textbf{Obstacle scene:} flat ground with four box-shaped obstacles
    of increasing height ($0.2$ to $1.3$\,m), which the robot runs into.
  \item \textbf{Terrain scene:} uneven, hilly terrain given as a triangle
    mesh.
\end{itemize}
Every system that PCR solves during $25$ time steps (one per Newton
iteration) is stored: $335$ systems with $n = 31$ to $67$ active constraint
rows in the obstacle scene, and $322$ systems with $n = 50$ to $170$ in the
terrain scene. Some sources of $\mathbf{Z}$ are built from earlier systems
of the sequence, so all results are reported on the last $20\%$ of each
sequence ($67$ and $65$ systems), which we call the \emph{held-out}
systems. Iteration counts are measured offline in double precision with
tolerance $10^{-5}$.

% Catalogue of the Z sources tried so far. Source names match
% Figure~\ref{fig:deflation-z-sources}. Median iterations (obstacle / terrain
% scene, Jacobi 19 / 19; data regenerated 2026-10-25, float64, tol 1e-5):
% exact bottom-16 11/9, cross-step 11/14, recycle-exact 11/9, harmonic m=16
% 20/11, m=24 16/12, m=32 11/10, GCRO-DR 14/10, joint/control 7/9,
% joint/control + 8 rows 5/4, null(J^T) 12/17 (k/n 0.59/0.82), enriched 4/6,
% hybrid k_c=8 7/5, random k=16 27.2 mean (obstacle, Jacobi 19.8).
\begin{table}[h]
  \centering
  \small
  \setlength{\extrarowheight}{2pt}
  \begin{tabular}{|>{\raggedright\arraybackslash}p{2.5cm}|>{\raggedright\arraybackslash}p{3.1cm}|>{\raggedright\arraybackslash}p{8.6cm}|}
    \hline
    Where $\mathbf{Z}$ comes from & $\mathbf{Z}$ source & What $\mathbf{Z}$ is \\
    \hline
    Eigensolve of the current system
      & exact bottom-$k$
      & The $k$ eigenvectors of the Jacobi-scaled $\mathbf{A}$ with the smallest
        eigenvalues: the directions in which PCR converges most slowly. This is
        the ideal $\mathbf{Z}$, but computing it costs more than the solve, so it
        only serves as a reference. \\
    \hline
    Exact eigenvectors of an earlier system
      & cross-step
      & The same $16$ eigenvectors, but computed once per time step, from the
        last system of the previous time step, and reused for every system of
        the current step. Works when $\mathbf{A}$ changes little from one step
        to the next. \\
    \cline{2-3}
      & recycle-exact
      & The $16$ eigenvectors of the system solved just before (the previous
        Newton iteration). It still needs an eigensolve for every system; it
        shows how much reusing earlier information can give at best. \\
    \hline
    Approximate eigenvectors from earlier PCR iterations
      & harmonic $m{=}16, 24, 32$
      & A solve of the previous system produces $m$ vectors (a Krylov basis)
        anyway. A small $m \times m$ eigenproblem in their span gives
        approximations of the $16$ smallest-eigenvalue eigenvectors (harmonic
        Ritz vectors). More vectors give better approximations, but the
        engine stops PCR after $16$ iterations. \\
    \cline{2-3}
      & GCRO-DR
      & Like harmonic, but the $16$ approximations are carried along the whole
        sequence of systems and refined with the new vectors of every solve,
        so they improve over time. The update rule comes from the GCRO-DR
        method. \\
    \hline
    Structure of the problem
      & joint/control
      & One unit vector (a column of the identity) per joint or control row.
        Deflating them solves the joint equations exactly, and PCR only
        handles the contact rows (Section~\ref{sec:z-joint-control}). \\
    \cline{2-3}
      & joint/control + 8 rows
      & The same, plus unit vectors on the $8$ contact rows that are most
        strongly coupled to another row, judged from the entries of
        $\mathbf{A}$ alone (Section~\ref{sec:z-strength}). \\
    \cline{2-3}
      & $\operatorname{null}(\mathbf{J}^\top)$
      & A basis of all force vectors $\mathbf{x}$ with
        $\mathbf{J}^\top\mathbf{x} = \mathbf{0}$: combinations of constraint
        forces that cancel out on every body. For them
        $\mathbf{A}\mathbf{x} = \mathbf{C}\mathbf{x}$, so only the small
        compliance resists them, which makes them soft. There are many of
        them: $k$ is $59\%$ to $82\%$ of $n$. \\
    \cline{2-3}
      & enriched
      & $\operatorname{null}(\mathbf{J}^\top)$ and joint/control together.
        Then $k \approx 0.9\,n$, so the coarse solve is almost a direct solve
        of the whole system. \\
    \hline
    Structure and eigensolve
      & joint/control + 8 eigenvectors
      & Joint/control plus the exact $8$ smallest-eigenvalue eigenvectors of
        the contact problem that remains after the joints are eliminated (the
        Schur complement $\mathbf{S}$). Needs an eigensolve; a reference for
        the cheap row rule above. \\
    \hline
    Sanity checks
      & random
      & $16$ random vectors. Deflation cannot make the spectrum worse, but
        this checks whether a useless $\mathbf{Z}$ can still cost iterations.
        In the obstacle scene it does. \\
    \cline{2-3}
      & tilted exact
      & The exact bottom-$16$ eigenvectors, deliberately tilted by a chosen
        angle $\theta$, to measure how accurate a $\mathbf{Z}$ has to be. \\
    \hline
  \end{tabular}
  \caption{Catalogue of the sources of $\mathbf{Z}$ we have tested.
    The names match Figure~\ref{fig:deflation-z-sources}.}
  \label{tab:z-sources-tried}
\end{table}

% NOTE: the figure and table use data regenerated 2026-10-25 on a GTX 1650.
% Obstacle reproduces deflation_ab.md; surface differs from it (Jacobi median 19 vs 29).
\begin{figure}[h]
  \centering
  \includegraphics[width=\textwidth]{latex-theory/figures/deflation_z_sources.pdf}
  \caption{Median PCR iterations on the held-out systems with deflation
    (DPCR) for real sources of $\mathbf{Z}$, against Jacobi PCR without
    deflation, which the ostrich solver uses today. The first four sources
    approximate the bottom-16 eigenvectors ($k = 16$). Joint/control uses
    unit vectors on the 18 joint and control rows, i.e.\ an exact solve of
    the joint block ($k = 18$). The last bar also eliminates the 8 contact
    rows with the largest AMG strength of connection
    $\max_{j \ne r} |a_{rj}| / \sqrt{a_{rr} a_{jj}}$ ($k = 26$), which needs
    only the entries of $\mathbf{A}$.
    Generated by \texttt{figure\_scripts/deflation\_z\_sources.py}.}
  \label{fig:deflation-z-sources}
\end{figure}

\begin{figure}[h]
  \centering
  \includegraphics[width=\textwidth]{latex-theory/figures/deflation_z_sources_time.pdf}
  \caption{GPU time per solve for the sources of Figure~\ref{fig:deflation-z-sources},
    for a single simulation ($W = 1$), median over the held-out systems.
    Blue: the solve itself, i.e.\ the setup of the deflation plus the
    iterations of Figure~\ref{fig:deflation-z-sources} times the measured
    cost of one iteration. Orange: building $\mathbf{Z}$ with dedicated GPU
    kernels. For cross-step and recycle-exact, the previous $\mathbf{Z}$ is
    refined by one or two steps of Chebyshev-filtered subspace iteration,
    the fewest steps that give the same PCR iteration counts as the exact
    eigenvectors; cross-step does this once per time step, recycle-exact
    once per system. Harmonic builds $\mathbf{Z}$ from the Krylov vectors of
    the previous solve; for $m = 16$ this needs no eigenproblem, for
    $m = 32$ a $32 \times 32$ one. Joint/control + 8 rows computes the
    strength-of-connection scores. Dotted line: Jacobi PCR, the current
    ostrich solver. Absolute times vary between runs on the laptop GPU used
    (up to about a factor of $2$ for the same quantity, as the GPU changes its
    clock); comparisons within one figure are consistent.
    Generated by \texttt{figure\_scripts/deflation\_z\_sources\_time.py}.}
  \label{fig:deflation-z-sources-time}
\end{figure}

\subsection{Joint and Control Rows: An Exact Block Solve}
\label{sec:z-joint-control}

In both scenes, $18$ rows of $\mathbf{A}$ belong to the joints ($15$)
and to the wheel controls ($3$). They are active in every system and sit at
fixed positions in the row ordering. Taking $\mathbf{Z}$ as the unit vectors
on these rows therefore costs nothing to build, and $k = 18$ is constant.
$\mathbf{A}\mathbf{Z}$ is just the corresponding columns of $\mathbf{A}$, and
$\mathbf{E} = \mathbf{Z}^\top\mathbf{A}\mathbf{Z} = \mathbf{A}_{JJ}$ is the
joint block.

With this $\mathbf{Z}$, deflation becomes block elimination. Order the joint
rows $J$ first and the contact rows $C$ second. Then
$\tilde{\mathbf{P}} = \mathbf{I} - \mathbf{A}\mathbf{Z}\mathbf{E}^{-1}\mathbf{Z}^\top$
gives
\begin{equation}
  \tilde{\mathbf{P}}\mathbf{A} =
  \begin{bmatrix} \mathbf{0} & \mathbf{0} \\ \mathbf{0} & \mathbf{S} \end{bmatrix},
  \qquad
  \mathbf{S} = \mathbf{A}_{CC} - \mathbf{A}_{CJ}\mathbf{A}_{JJ}^{-1}\mathbf{A}_{JC} .
\end{equation}
The joint block is solved exactly, and PCR iterates only on the Schur
complement $\mathbf{S}$ of the contact rows. In mechanics this is known as
static condensation.

The joint rows are a good target for an exact solve for three reasons:
\begin{itemize}
  \item \textbf{Jacobi cannot handle them.} Jacobi treats every row on its
    own and ignores the off-diagonal coupling. Joint rows are stiff
    (compliance $c = 0.011$, controls $0.22$) and strongly coupled to each
    other, because several rows of one joint act on the same pair of
    bodies. In the Jacobi-scaled matrix, the largest off-diagonal entry of a
    joint row is typically $0.54$ in the obstacle scene and $0.30$ in the
    terrain scene. For a
    typical contact row it is close to $0$.
  \item \textbf{The block itself is easy.} The Jacobi-scaled
    $\mathbf{A}_{JJ}$ has condition number about $2.5$. Its $18 \times 18$
    Cholesky factorization is cheap and stable, even in single precision.
  \item \textbf{In the terrain scene, the soft modes live there.} The joint rows are
    $14\%$ of the rows but carry $51\%$ of the weight of the bottom-16
    eigenvectors of the scaled $\mathbf{A}$. Eliminating them lowers the
    condition number from about $20$ to about $4$. In the obstacle scene the soft modes
    are not concentrated on the joint rows, and the condition number only
    drops from $10^6$ to $3 \cdot 10^5$. The iteration count nevertheless
    drops from $19$ to $7$ (Figure~\ref{fig:deflation-z-sources}), a
    reminder that the condition number alone does not set the iteration
    count (Section~\ref{sec:dpcr}).
\end{itemize}
The exact elimination is what matters, not merely using $\mathbf{A}_{JJ}$.
Block-Jacobi with the exact $\mathbf{A}_{JJ}^{-1}$, i.e.\ the same block used
as a preconditioner instead of being eliminated, needs about twice as many
iterations in the obstacle scene ($14$).

\subsection{Choosing Further Rows by Strength of Connection}
\label{sec:z-strength}

The joint block removes the coupling that Jacobi misses among the joint
rows. Some contact rows are strongly coupled too, and the same idea extends
to them: eliminate a few more rows exactly. The question is which ones,
without computing any eigenvectors.

The answer comes from looking at the matrix after Jacobi scaling,
$\mathbf{A}_s = \mathbf{P}^{-1/2}\mathbf{A}\mathbf{P}^{-1/2}$, where
$\mathbf{P} = \operatorname{diag}(\mathbf{A})$ is the Jacobi preconditioner. It has a unit diagonal and
off-diagonal entries
\begin{equation}
  s_{rj} = \frac{|a_{rj}|}{\sqrt{a_{rr}\,a_{jj}}} \in [0, 1),
  \label{eq:strength}
\end{equation}
where $s_{rj} < 1$ because $\mathbf{A}$ is SPD. Jacobi would be exact if every
$s_{rj}$ were zero. Now consider two rows $r$ and $j$ alone. Their scaled
$2 \times 2$ block $\begin{bmatrix} 1 & s_{rj} \\ s_{rj} & 1 \end{bmatrix}$
has eigenvalues $1 \pm s_{rj}$. If $s_{rj}$ is close to $1$, the two rows
are almost parallel, and $1 - s_{rj}$ is a soft mode that Jacobi cannot
remove. This is exactly the example of Section~\ref{sec:deflated-cg}, where
$s = 0.99/1.01 \approx 0.98$ gives the scaled eigenvalues $1.98$ and $0.02$.
Eliminating both rows exactly removes this mode. Eliminating only one of them
is not enough: the other row stays in the Schur complement with diagonal
$a_{rr}(1 - s_{rj}^2)$, and since DPCR keeps the preconditioner
$\mathbf{P} = \operatorname{diag}(\mathbf{A})$, it still sees the small
eigenvalue $1 - s_{rj}^2 \approx 2(1 - s_{rj})$.

The quantity \eqref{eq:strength} is the \emph{strength of connection} used
by algebraic multigrid to decide which unknowns are strongly coupled; this
symmetric, scaled form is the one used in smoothed aggregation
\citep{vanek1996algebraic}. Our rule scores each contact row by its strongest
connection to any other row, joint rows included,
\begin{equation}
  \sigma_r = \max_{j \ne r} s_{rj},
\end{equation}
and adds unit vectors on the $k_c$ contact rows with the largest $\sigma_r$
to the joint and control rows. With $k_c = 8$ this gives $k = 26$. Both
rows of a strongly coupled pair have the same score, so the rule tends to
pick them together, as the $2 \times 2$ argument requires. The
rule needs only the diagonal and the off-diagonal entries of $\mathbf{A}$,
i.e.\ one pass over the matrix, with no solves and no eigenvectors.
